% ============================================================
% f2014_sichuan_science_q13_river_width.tex
% 2014 四川卷（理科）第 13 题：气球上测河流宽度
% 来源 PDF：普通高考/2014/2014四川理.pdf
% 原页：第 1 页（A3 横排）中栏下部，印面 919，第 13 题题干下方
% 说明：A 为气球位置，其正下方为河岸上的垂足；B、C 为近岸与对岸
%       观测点，A 到 B、C 的俯角分别为 67°、30°；气球高 46 m；
%       图中给出过 A 的水平虚线、铅垂高度线与垂足处的直角标记。
% ============================================================
\begin{tikzpicture}[>=Stealth, line width=0.7pt, font=\small,
  x=1cm, y=1cm, line cap=round, line join=round]

  \coordinate (A) at (0,4);
  \coordinate (D) at (0,0);
  \coordinate (B) at (1.70,0);
  \coordinate (C) at (6.93,0);

  % 地面与过 A 的水平线
  \draw[dashed] (-0.80,0) -- (7.60,0);
  \draw[dashed] (-0.20,4) -- (2.60,4);

  % 铅垂线与视线
  \draw[dashed] (A) -- (D);
  \draw[dashed] (A) -- (B);
  \draw[dashed] (A) -- (C);

  % 气球
  \draw (0,4.46) circle (0.22);
  \draw (-0.08,4.27) -- (0.08,4.27) -- (0.10,4.00) -- (-0.10,4.00) -- cycle;
  \node[font=\footnotesize] at (0,4.46) {$A$};

  % 铅垂高度与标注
  \draw (-0.85,4) -- (-0.25,4);
  \draw (-0.85,0) -- (-0.25,0);
  \draw[<->] (-0.55,4) -- (-0.55,0);
  \node[left] at (-0.70,2) {$46\mathrm{m}$};

  % 垂足处的直角标记
  \draw (0,0.16) -- (0.16,0.16) -- (0.16,0);

  % 俯角：30° 到远岸 C，67° 到近岸 B
  \draw ($(A)+(0.90,0)$) arc (0:-30:0.90);
  \node at (1.22,3.74) {$30^\circ$};
  \draw ($(A)+(1.70,0)$) arc (0:-67:1.70);
  \node at (2.05,3.68) {$67^\circ$};

  % 点的标注
  \node[below] at (B) {$B$};
  \node[below] at (C) {$C$};
\end{tikzpicture}
